\section{Conclusion}

We formulated continual learning as an old-solution preservation constraint (CLOC) and derived two local quadratic relaxations. Preserving each old sample's scalar loss yields an outer-product loss-gradient matrix; for negative log likelihood, this is algebraically the observed-label empirical Fisher used in many EWC implementations, but it is not the model Fisher or a general curvature identity. Preserving outputs modulo the tangent to the old loss level set instead yields the iEF normalization. This retains the central geometric intuition while locating its exact boundary: undamped normalization removes the samplewise output-loss-gradient scale, whereas positive damping multiplies each summand by $s/(s+\tau)$ and assigns zero contribution at exact interpolation.

Task-relative damping restores invariance to global changes of loss units, and the controlled contamination study directly shows that normalization can prevent high-loss corrupted samples from dominating a full importance matrix in the tested construction. Whether this altered local geometry improves a diagonal continual-learning penalty remains an empirical, domain-conditional question. Our evaluation therefore separates normalization from residual-diversity weighting, matches global scale, includes external regularization and offline references, selects hyperparameters only on validation data, and reserves paired task orders and test splits for locked inference. IEWC is best understood as a CLOC-derived family whose damping, weighting, output metric, and matrix approximation are explicit choices---not as a universal replacement for every continual-learning method.
